\section{Respuestas directas a las preguntas de la profesora}

\begin{description}[leftmargin=1.2em,style=nextline,itemsep=0.35em]
  \item[1.a) Tasas, estabilidad y justificación del M/M/1.]
  $\lambda=0{,}25$ trabajos/min y $\mu=1/3$ trabajos/min, en minutos. Entonces
  $\rho=\lambda/\mu=0{,}75<1$: el sistema es estable. Es M/M/1 porque las
  llegadas son Poisson, los servicios exponenciales, hay un solo servidor, la
  disciplina es FIFO y la cola es prácticamente ilimitada.

  \item[1.b) Simulador de eventos discretos.]
  Implementado con un motor propio basado en \texttt{heapq}, sin SimPy, con
  reloj de eventos, flujos aleatorios independientes y semillas explícitas.

  \item[1.c) Métricas del M/M/1 tras el warm-up.]
  Con 50 réplicas, warm-up de 1\,000 min y horizonte de 30\,000 min:
  $L=3{,}0466$, $\Lq=2{,}2961$, $W=12{,}2045$ min y $\Wq=9{,}1973$ min.

  \item[1.d) Ley de Little y comparación analítica.]
  Los cocientes de Little son $0{,}999997$ y $0{,}999978$ (prácticamente 1).
  Los cinco IC95\,\% contienen los valores analíticos
  ($L=3$, $\Lq=2{,}25$, $W=12$, $\Wq=9$, $\rho=0{,}75$).

  \item[2. Modelo del SOC.]
  $s=3$ analistas, $K=8$ espacios de cola (capacidad total $11$), dos colas
  FIFO por tipo, prioridad no expropiativa de las críticas, llegadas Poisson
  independientes ($\lambda_1=0{,}05$, $\lambda_2=0{,}10$), servicio uniforme
  $[15,25]$ para Tipo~1 y Weibull $(k=1{,}5,\beta=12)$ para Tipo~2. Las
  críticas también pueden bloquearse si la cola está llena.

  \item[2. Métricas del SOC a 30 días.]
  $\Wone=22{,}260$ min, $\Wqone=2{,}267$ min, $\Wtwo=17{,}250$ min,
  $\Wqtwo=6{,}417$ min, $\Lq=0{,}750$, utilización conjunta $68{,}87$\,\% y
  $\loss=0{,}385$\,\%, calculadas solo sobre alertas admitidas.

  \item[Escenario A.]
  Al aumentar $\lambda_2$ un $40$\,\%, $\loss$ sube de $0{,}385$\,\% a
  $2{,}220$\,\%: $+1{,}835$ pp, un aumento relativo del $477{,}0$\,\%. Los
  analistas no se saturan teóricamente ($\rho_{\text{ofrecida}}=0{,}8389$,
  utilización $81{,}65$\,\%), pero la congestión y la pérdida de alertas
  críticas aumentan de forma severa.

  \item[Escenario B.]
  Bajo demanda base ambas alternativas reducen $\loss$ en más del $94$\,\% y
  empatan estadísticamente. Bajo sobrecarga, $s=4$ reduce $\loss$ a
  $0{,}129$\,\% frente a $0{,}395$\,\% con $K=16$. El cuarto analista también
  minimiza las esperas en ambos regímenes. Recomendación: contratar el cuarto
  analista y mantener $K=8$.
\end{description}
